% ============================================================
% PHẦN SADGS — Slide 4/8: Giãn Gaussian dưới cỡ (expand_undersized_gs)
% ============================================================

\begin{frame}{Vì sao cần "giãn ngược" -- Gaussian quá nhỏ cũng lãng phí ngân sách \alert{(cơ chế này KHÔNG chạy mặc định, xem Slide 4.5)}}
\footnotesize
\begin{itemize}
    \item \alert{Lưu ý ngay từ đầu:} \texttt{expand\_undersized\_gs} được \textbf{định nghĩa đầy đủ} trong \texttt{gaussian\_model.py}, nhưng lời gọi duy nhất tới nó trong \texttt{train.py:356--359} \textbf{đang bị comment out} -- mặc định SADGS \textbf{không} chạy bước "giãn ngược" mô tả dưới đây, chỉ clone/split. Chi tiết ở frame cuối và ở \texttt{part\_sadgsx\_06}.
    \item ADC gốc (3DGS) chỉ có 2 phản ứng khi Gaussian sai cỡ: \textbf{clone} (nhân bản, cho vùng thiếu điểm) và \textbf{split} (chia nhỏ, cho vùng thừa điểm) -- không có cơ chế \textbf{thu gọn số lượng} khi một vùng đã bị chia \emph{quá mức cần thiết}.
    \item SADGS đo $\eta = \text{axis}_k / w_{\min}$ (\texttt{freq\_utils.py:348}): tỉ số giữa độ dài trục chiếu màn hình của Gaussian và bước sóng cục bộ của texture -- đã dùng để \textbf{chia} khi $\eta$ quá lớn (Slide 3). \alert{Lưu ý:} docstring của \texttt{expand\_undersized\_gs} (\texttt{gaussian\_model.py:838}) lại giả định $\eta=(\sigma\omega)^2$ -- một mâu thuẫn với $\eta$ tuyến tính thực sự được tính (hệ quả ở slide sau).
    \item Chiều ngược lại: nếu $\eta$ \textbf{quá nhỏ} ($0 < \eta < \tau_{\text{expand}}$, tham số \texttt{tau\_expand}, mặc định $1.0$), Gaussian đang \emph{mờ hơn mức cần thiết} so với chi tiết cục bộ -- một Gaussian to hơn có thể phủ đúng vùng đó mà \textbf{không mất chi tiết} và không tốn thêm 1 điểm nào trong ngân sách.
    \item Đây là đòn bẩy thứ ba của SADGS bên cạnh clone/split: \textbf{sửa scale tại chỗ bằng công thức giải tích}, không cần thêm/xoá Gaussian, không cần lặp gradient descent.
\end{itemize}
\end{frame}

% --- Slide mới: công thức giải tích ---
\begin{frame}{Công thức giải tích: co/giãn đúng \emph{một bước} tới biên Nyquist}
\footnotesize
\begin{itemize}
    \item Mục tiêu đặt ra trong docstring: $\eta_{\text{new}} = 1$. Với \emph{giả định} $\eta\propto\sigma^2$ của docstring (\texttt{gaussian\_model.py:837--840}):
\end{itemize}
\[
\eta_{\text{new}} = \Big(\frac{\sigma_{\text{target}}}{\sigma_{\text{old}}}\Big)^2 \eta_{\text{old}} = 1
\;\Longrightarrow\;
\sigma_{\text{target}} = \frac{\sigma_{\text{old}}}{\sqrt{\eta_{\text{old}}}}
\]
\begin{itemize}
    \item Vì scale được lưu trong \textbf{log-space} (\texttt{self.\_scaling}, qua \texttt{scaling\_inverse\_activation}), phép chia trên trở thành một \textbf{phép cộng trực tiếp}:
\end{itemize}
\[
\log\sigma_{\text{target}} = \log\sigma_{\text{old}} - \tfrac12\log\eta_{\text{old}}
\quad\Longleftrightarrow\quad
\Delta\!\log\sigma = -\tfrac12\log\eta
\]
\begin{itemize}
    \item Code (\texttt{expand\_undersized\_gs}, \texttt{gaussian\_model.py:833--864}): \texttt{eta\_vals = clamp(eta, min=1e-6)} (tránh $\log 0$), rồi \texttt{delta\_log\_scale = -0.5*torch.log(eta\_vals)} cộng thẳng vào \texttt{self.\_scaling} -- \textbf{không lặp, không tối ưu}, 1 phép gán tensor là xong.
    \item \alert{Nhưng} $\eta$ thực tế \textbf{tuyến tính} theo $\sigma$ ($\eta=\text{axis}/w_{\min}$), nên $\sigma\!\to\!\sigma/\sqrt{\eta}$ chỉ cho $\eta_{\text{new}}=\sqrt{\eta}$ --- \emph{rút căn} khoảng cách tới $1$, không đạt $1$ như comment \texttt{":853"} tuyên bố. Muốn đúng một bước phải dùng $\Delta\!\log\sigma=-\log\eta$.
\end{itemize}
\end{frame}

% --- Slide mới: dị hướng theo trục + so sánh với split ---
\begin{frame}{Cũng dị hướng theo từng trục -- nhưng ngược chiều với split}
\footnotesize
\begin{itemize}
    \item Giống split (Slide 3), \texttt{expand\_undersized\_gs} nhận \texttt{max\_eta\_3ch} -- $\eta$ tách riêng theo 3 trục $x,y,z$ -- và chỉ sửa trục nào thoả điều kiện:
\end{itemize}
\[
\text{undersized\_mask}_{\text{axis}} = \big(0 < \eta_{\text{axis}} < \tau_{\text{expand}}\big)
\]
\begin{itemize}
    \item Trục không thoả điều kiện (đã đúng cỡ hoặc đang quá lớn, chờ split) được \textbf{giữ nguyên} -- \texttt{delta} chỉ khác 0 tại các trục dưới cỡ, nên một Gaussian có thể vừa "đủ cỡ" theo trục này vừa "được giãn" theo trục khác trong cùng 1 bước.
    \item Khác biệt cốt lõi với split: split dùng $k=\lceil\sqrt{\eta}\rceil$ (làm tròn lên số nguyên $\ge 1$) để nhân bản thành lưới $k_x\times k_y\times k_z$ Gaussian con; expand dùng $\sqrt{\eta}$ \textbf{liên tục} (không làm tròn) để chia scale ngay trên chính Gaussian đó -- \textbf{số lượng Gaussian không đổi}.
\end{itemize}
\centering
\includegraphics[width=0.3\linewidth]{fig_09_split_geometry.png}
\captionof{figure}{\footnotesize Đối chiếu: hình mượn từ cơ chế split dị hướng (Slide 3) -- nơi $\eta$ lớn sinh lưới $k_x\times k_y\times k_z$ Gaussian con; expand xử lý chiều ngược lại ($\eta$ nhỏ) bằng cách co scale tại chỗ, không sinh thêm điểm}
\end{frame}

% --- Slide mới: ví dụ số ---
\begin{frame}{Ví dụ số: một Gaussian dưới cỡ trên 2/3 trục}
\footnotesize
\begin{itemize}
    \item Giả sử $\tau_{\text{expand}}=1.0$ và một Gaussian có $\eta = (\eta_x,\eta_y,\eta_z) = (0.25,\ 0.64,\ 1.5)$:
\end{itemize}
\centering
\begin{tabular}{lccc}
\toprule
\textbf{Trục} & $\eta_{\text{axis}}$ & Dưới cỡ? ($0<\eta<1$) & $\Delta\log\sigma = -\tfrac12\log\eta$ \\
\midrule
$x$ & $0.25$ & Có & $-\tfrac12\log(0.25) \approx +0.693$ \\
$y$ & $0.64$ & Có & $-\tfrac12\log(0.64) \approx +0.223$ \\
$z$ & $1.5$  & Không (chờ split) & $0$ (giữ nguyên) \\
\bottomrule
\end{tabular}
\vspace{0.2cm}
\footnotesize
\begin{itemize}
    \item Kết quả: $\sigma_x$ tăng theo hệ số $1/\sqrt{0.25}=2\times$, $\sigma_y$ tăng $1/\sqrt{0.64}\approx1.25\times$, $\sigma_z$ giữ nguyên -- Gaussian trở nên "phình" đúng theo hướng đang quá mờ so với chi tiết, còn trục $z$ để dành cho vòng split kế tiếp.
    \item Số lượng Gaussian: vẫn là \textbf{1} (so với split cùng $\eta$ thì $k_z=\lceil\sqrt{1.5}\rceil=2$ nếu $\eta_z$ vượt ngưỡng split, sinh thêm điểm) -- minh hoạ đúng vai trò "tiết kiệm ngân sách" của expand.
\end{itemize}
\centering
\includegraphics[width=0.32\linewidth]{fig_11_N_accounting.png}
\captionof{figure}{\footnotesize Kế toán $N$ theo cơ chế split (tham khảo đối chiếu): expand không cộng thêm vào $N$, chỉ cập nhật $\sigma$ tại chỗ}
\end{frame}

% --- Slide mới: bảng so sánh + vị trí trong pipeline ---
\begin{frame}{So sánh expand vs.\ split và vị trí trong vòng lặp ADC}
\footnotesize
\begin{tabular}{lll}
\toprule
\textbf{Tiêu chí} & \texttt{expand\_undersized\_gs} & \texttt{densify\_and\_split\_structgs} \\
\midrule
Điều kiện kích hoạt & $0<\eta_{\text{axis}}<\tau_{\text{expand}}{=}1.0$ & $\eta_{\text{axis}}\ge 1$ (qua $k\ge\sqrt{\eta}$) \\
Số Gaussian & Không đổi (sửa tại chỗ) & Tăng: $N{=}k_x k_y k_z$ mới thay 1 cũ \\
Công thức scale & $\sigma_{\text{old}}/\sqrt{\eta}$ (liên tục $\Rightarrow \eta_{\text{new}}{=}\sqrt{\eta}$) & $\sigma_{\text{old}}/k^{p}$, $k{=}\lceil\sqrt{\eta}\rceil$, $p{=}$\texttt{ks\_scale\_power}${=}1$ \\
Vị trí/hướng & Không đổi & Dịch theo lưới $k$, xoay bởi $R(q)$ \\
Chi phí & $O(1)$, 1 phép cộng log-space & $O(N)$ điểm mới + prune điểm cũ \\
\bottomrule
\end{tabular}
\vspace{0.3cm}
\begin{itemize}
    \item \alert{Thực trạng:} lời gọi \texttt{gaussians.expand\_undersized\_gs(...)} trong \texttt{train.py:356--359} \textbf{đang bị comment out}, nên cơ chế này \textbf{không chạy} trong luồng huấn luyện mặc định. Về thiết kế, nó được đặt \textbf{ngay trước} clone/split trong cùng bước densify.
    \item Vì cả hai cơ chế đều dựa trên cùng một đại lượng $\eta$ 3 kênh, SADGS xử lý \textbf{cả hai đầu} của sai số mật độ (quá thưa lẫn quá thừa) bằng công cụ toán học nhất quán, thay vì chỉ có "chia đôi khi gradient cao" như ADC 3DGS gốc.
\end{itemize}
\end{frame}
